\printedpage{}
% head: (none); no page number printed on this page

\begin{center}
\vspace*{6cm}

{\LARGE\textsf{THE DERIVATORS}}

\vspace{1.5cm}

{\large\textbf{ALEXANDRE GROTHENDIECK}}

\vspace{2.5cm}

{\large\textbf{Chapter V}}

\vspace{0.4cm}

{\large\textbf{Categories of paths and localization}}
\end{center}

\clearpage

\printedpage{}
% head: (none); no page number printed on this page

\vspace*{6cm}

This text was deciphered and transcribed in \LaTeXe{} by M. Künzer. It was edited by M. Künzer, J. Malgoire, and G. Maltsiniotis. The transcription is as faithful as possible to the manuscript. For the few obvious corrections, or rare comments by the editors, as well as for the original numbering of the manuscript pages, the typewriter characters \texttt{[typewriter]} in brackets are used. A question mark in brackets means that one is not sure of the preceding word.

This edition is provisional. Remarks, comments and corrections are welcome. Send a message to:

\textbf{[e-mail address removed]}\\
\hspace*{\fill}G. Maltsiniotis\hspace*{2cm}

\clearpage

\printedpage{1}
% head: \textbf{A. Grothendieck :} \textit{Dérivateurs,} Ch. V, Édition des Universités de Montpellier II et Paris VII
% foot: Version 9 septembre 2001

\msp{page 1}

\section*{1\quad The category of paths}

Let $\Phi$ be the following category:
\[
\left\{
\begin{array}{rcll}
\Ob\Phi & = & \mathbf{Z} & \\
\Fl\Phi & = & \Ob\Phi\sqcup\mathbf{Z} & \qquad (\Ob\Phi\leftrightarrow\text{identities})
\end{array}
\right.
\]
We denote by $\{\alpha\}$, or simply $\alpha$, the object corresponding to $\alpha\in\mathbf{Z}$. We denote by $u_\alpha$ the arrow corresponding to the integer $\alpha\in\mathbf{Z}$.
\[
\begin{array}{rcl@{\qquad}l@{\quad}l@{\quad}l}
s(u_\alpha) & = & \alpha & \text{if } \alpha \text{ even}, & \alpha+1 & \text{if } \alpha \text{ odd}\\[1ex]
b(u_\alpha) & = & \alpha+1 & \text{if } \alpha \text{ even}, & \alpha-1 & \text{if } \alpha \text{ odd}
\end{array}
\]
\[
\texttt{[rather}\qquad b(u_\alpha) \;=\; \alpha+1 \quad \texttt{if}\ \alpha\ \texttt{even,}\quad \alpha \quad \texttt{if}\ \alpha\ \texttt{odd]}
\]

No compositions to define, for if two arrows are composable, they are the identity. In fact, $\Phi$ is an \textit{ordered} category, $\{\alpha\}<\{\beta\}$ if and only if $\alpha$ even, $\beta=\alpha+1$ \texttt{[rather} $\beta=\alpha+1$ \texttt{or} $\beta=\alpha-1$\texttt{]}. % [sic] Il n'y pas -> Il n'y a pas, verb "a" missing

We denote, for $n\in\mathbf{N}$, by
\[
\Phi_n\subseteq\Phi,\quad \Ob\Phi_n=\{\alpha\in\Ob\Phi\mid -n\le\alpha\le+n\},
\]
the full subcategory of $\Phi$ whose objects are the $\{\alpha\}$ with $-n\le\alpha\le+n$. Thus
\[
\Phi_1=\{-1\longleftarrow 0\longrightarrow 1\},
\]
the dual is
\[
\Psi_1=\Phi_1^{o}=\{-1\longrightarrow 0\longleftarrow 1\}.
\]
We can regard $\Phi$ as an amalgamated sum of an infinite ordered sequence (on both sides) of copies of $\Phi_1$, or also of a similar sequence of copies of $\Psi_1$.
\[
\begin{tikzcd}[column sep=small, execute at end picture={
\draw[decorate,decoration={brace,amplitude=4pt}] ([yshift=10pt]\tikzcdmatrixname-1-2.center) -- ([yshift=10pt]\tikzcdmatrixname-1-4.center) node[midway,above=4pt]{$\scriptstyle\Phi_1$};
\draw[decorate,decoration={brace,amplitude=4pt}] ([yshift=10pt]\tikzcdmatrixname-1-4.center) -- ([yshift=10pt]\tikzcdmatrixname-1-6.center) node[midway,above=4pt]{$\scriptstyle\Phi_1$};
\draw[decorate,decoration={brace,amplitude=4pt}] ([yshift=10pt]\tikzcdmatrixname-1-6.center) -- ([yshift=10pt]\tikzcdmatrixname-1-8.center) node[midway,above=4pt]{$\scriptstyle\Phi_1$};
\draw[decorate,decoration={brace,amplitude=4pt}] ([yshift=10pt]\tikzcdmatrixname-1-8.center) -- ([yshift=10pt]\tikzcdmatrixname-1-10.center) node[midway,above=4pt]{$\scriptstyle\Phi_1$};
\draw[decorate,decoration={brace,amplitude=4pt}] ([yshift=10pt]\tikzcdmatrixname-1-10.center) -- ([yshift=10pt]\tikzcdmatrixname-1-12.center) node[midway,above=4pt]{$\scriptstyle\Phi_1$};
\draw[decorate,decoration={brace,mirror,amplitude=4pt}] ([yshift=-10pt]\tikzcdmatrixname-1-3.center) -- ([yshift=-10pt]\tikzcdmatrixname-1-5.center) node[midway,below=4pt]{$\scriptstyle\Psi_1$};
\draw[decorate,decoration={brace,mirror,amplitude=4pt}] ([yshift=-10pt]\tikzcdmatrixname-1-5.center) -- ([yshift=-10pt]\tikzcdmatrixname-1-7.center) node[midway,below=4pt]{$\scriptstyle\Psi_1$};
\draw[decorate,decoration={brace,mirror,amplitude=4pt}] ([yshift=-10pt]\tikzcdmatrixname-1-7.center) -- ([yshift=-10pt]\tikzcdmatrixname-1-9.center) node[midway,below=4pt]{$\scriptstyle\Psi_1$};
\draw[decorate,decoration={brace,mirror,amplitude=4pt}] ([yshift=-10pt]\tikzcdmatrixname-1-9.center) -- ([yshift=-10pt]\tikzcdmatrixname-1-11.center) node[midway,below=4pt]{$\scriptstyle\Psi_1$};
}]
\cdots & {} & -4 \arrow[l] \arrow[r] & -3 & -2 \arrow[l] \arrow[r] & -1 & 0 \arrow[l] \arrow[r] & 1 & 2 \arrow[l] \arrow[r] & 3 & 4 \arrow[l] \arrow[r] & {} & \cdots
\end{tikzcd}
\]

The $\Phi_n$, for $n$ odd, are likewise sums of copies of $\Phi_1$ put end to end, and the $\Psi_n$, for $n$ even $>0$, sums of copies of $\Psi_1$.

\clearpage

\printedpage{2}
% head: \textbf{A. Grothendieck :} \textit{Dérivateurs,} Ch. V, Édition des Universités de Montpellier II et Paris VII
% foot: Version 9 septembre 2001

The $\Phi_n$ form a direct system with fully faithful transition functors, of which $\Phi$ is the direct limit
\[
\Phi=\varinjlim_n\Phi_n.
\]
But for what follows we are going to make the $\Phi_n$ into an \textit{inverse} system, by defining for every $n\le m$ a \textit{retraction} $\varphi_{n,m}$ of $\Phi_m$ onto $\Phi_n$, by

\msp{page 2}
\[
\left\{
\begin{array}{l}
\varphi_{n,m}:\Phi_m\longrightarrow\Phi_n\subset\Phi_m\\[1ex]
\varphi_{n,m}(\{\alpha_i\})=
\left\{
\begin{array}{rcll@{\quad}l@{\ \ }r}
\alpha_i & = & i & \in\Ob\Phi_n, & \text{i.e.} & -n\le i\le n\\[1ex]
\alpha_i & = & n & & \text{if} & n\le i\le m\\[1ex]
\alpha_i & = & -n & & \text{if} & -m\le i\le -n
\end{array}
\right.
\end{array}
\right.
\]

\texttt{[rather}
\[
\begin{array}{rcl}
\varphi_{n,m}:\Phi_m & \longrightarrow & \Phi_n\subset\Phi_m\\[1ex]
\varphi_{n,m}(\{i\}) & = &
\left\{
\begin{array}{l@{\qquad}l@{\quad}l@{\ \ }r}
\{i\} & \texttt{if}\ i\in\Ob\Phi_n, & \texttt{i.e} & -n\le i\le n\\[1ex]
\{n\} & & \texttt{if} & n\le i\le m\\[1ex]
\{-n\} & & \texttt{if} & -m\le i\le -n\ \texttt{]}
\end{array}
\right.
\end{array}
\]
We denote the inverse system of ordered sets thus obtained by
\[
\bPhi=(\Phi_n)_{n\in\mathbf{N}}\in\mathrm{Pro}(\mathrm{Cat}).
\]
Let $X$ be a category, then for $n\in\mathbf{N}$
\[
\uCh_n(X)\overset{\text{def}}{=}\uHom(\Phi_n,X)
\]
is the category of diagrams in $X$ of the form
\[
A_{-n}\cdots\longleftarrow A_{-2}\xrightarrow{u_{-2}}A_{-1}\xleftarrow{u_{-1}}A_0\xrightarrow{u_0}A_1\xleftarrow{u_1}A_2\xrightarrow{u_2}A_3\longleftarrow\cdots A_n
\tag*{$(*)$}
\]
with $2n$ arrows end to end $u_i$ $(-n\le i\le n)$ \texttt{[rather (}$-n\le i\le n-1$\texttt{)]}, two consecutive arrows always going in opposite directions,
\[
\text{either}\ \ \cdot\longrightarrow\cdot\longleftarrow\cdot,\ \ \text{or}\ \ \cdot\longleftarrow\cdot\longrightarrow\cdot\,.
\]
The canonical functor deduced from $\varphi_{n,m}:\Phi_m\longrightarrow\Phi_n$
\[
\varphi_X^{m,n}:\uCh_n(X)\longrightarrow\uCh_m(Y),\qquad n\le m,
\]
\printedpage{3}
% head: \textbf{A. Grothendieck :} \textit{Dérivateurs,} Ch. V, Édition des Universités de Montpellier II et Paris VII
% foot: Version 9 septembre 2001
\texttt{[rather} $\uCh_m(X)$\texttt{]} is the ``degeneracy functor'', consisting in completing a diagram $(*)$, with $m-n$ terms $A_{-n}$ with identities on the left, and with $m-n$ terms \texttt{[}$A_n$\texttt{]} with identities on the right, so for $n$ even we find
\[
\underbrace{A_{-n}\cdots\xrightarrow{\text{id}}A_{-n}}_{m-n\text{ terms}}\xleftarrow{\text{id}}A_{-n}\longrightarrow A_{-(n-1)}\cdots\cdots A_{n-1}\longleftarrow A_n\xrightarrow{\text{id}}\underbrace{A_n\xleftarrow{\text{id}}A_n\longrightarrow\cdots A_n}_{m-n\text{ terms}}.
\]
This is a fully faithful functor (which corresponds to the fact that $\varphi_{n,m}$ is surjective on objects and on arrows. % [sic] flèches. -> flèches)., opening parenthesis never closed

\msp{page 3}
We set
\[
\underbrace{\uCh_\infty(X)}_{\text{or simply }\uCh(X)}\qquad=\uHom(\bPhi,X)\overset{\text{def}}{=}\varinjlim_n\uCh_n(X)
\]
(\footnote{The set of objects of this category will be denoted by $\mathrm{Ch}_\infty(X)$ or $\mathrm{Ch}(X)$, likewise we define the $\mathrm{Ch}_n(X)$.}) (\footnote{in this formula: $\bPhi$ procategory}). We have a fully faithful inclusion
\[
\uCh_\infty(X)\lhook\joinrel\longrightarrow\uHom(\Phi,X),
\]
where the second category is that of infinite diagrams (in both directions)
\[
\cdots\longleftarrow A_{-2}\xrightarrow{u_{-2}}A_{-1}\xleftarrow{u_{-1}}A_0\xrightarrow{u_0}A_1\xleftarrow{u_1}A_2\xrightarrow{u_2}\cdots
\]
and where $\uCh_\infty(X)$ is identified with the full subcategory formed by the objects for which the arrows $u_i$ are identities, for $|i|$ sufficiently large. % [sic] formées -> formée, agreement with "la sous-catégorie"

We shall call $\uCh_\infty(X)$ the \textit{category of paths} in $X$, its objects will be the paths in $X$. Those which come from $\uCh_n(X)$ will be called \textit{standard paths of length} $2n$. (The term ``standard'' applies here to our notational conventions. Thus we consider only paths of even length, and for a given length, we choose one of the two reasonable kinds, for example for paths of length 2 we have taken only
\[
A_{-1}\longleftarrow A_0\longrightarrow A_1
\]
and not
\[
A_{-1}\longrightarrow A_0\longleftarrow A_1.)
\]

\msp{page 4}
I realize that my notations $u_i$ are not convenient for duality, I am going to change them, to $u_i$ for the arrows between objects $\{i\}$ with $i\ge0$, and $v_i$ the arrow symmetric to $u_i$ with respect to the origin. (The group of automorphisms of the $\Phi_n$ is reduced to the group with two elements, the only nontrivial automorphism being the symmetry $\{i\}\longrightarrow\{-i\}$. So our paths will henceforth be denoted % [sic] {-i}. -> {-i}). , opening parenthesis "(Le groupe" not closed on this page
\[
A_{-2}\xrightarrow{v_1}A_{-1}\xleftarrow{v_0}A_0\xrightarrow{u_0}A_1\xleftarrow{u_1}A_2\longrightarrow\cdots.
\tag*{$(*)$}
\]